\documentclass[aspectratio=169, 10pt]{beamer}
\usetheme{Madrid}
\usecolortheme{beaver}

\usepackage{amsmath, amssymb, amsthm}
\usepackage{booktabs}
\usepackage{graphicx}
\usepackage{tikz}
\usepackage{algorithm2e}
\usepackage{hyperref}

\title[TabDiff Reproduction \& Extension]{TabDiff: Mixed-Type Diffusion for Tabular Data}
\subtitle{Paper Reproduction, Algorithmic Extensions, and Experimental Evaluation}
\author[DFMGAI Team]{Course Project: Diffusion \& Flow Matching Generative AI (DFMGAI)}
\institute[University]{Department of Computer Science \& Engineering}
\date{\today}

\begin{document}

% Title slide
\begin{frame}
  \titlepage
\end{frame}

% Outline / Agenda
\begin{frame}{Presentation Outline}
  \tableofcontents
\end{frame}

% =========================================================================
% (a) Problem statement and motivation
% =========================================================================
\section{(a) Problem Statement and Motivation}

\begin{frame}{(a) Problem Statement and Motivation}
  \textbf{Background:}
  \begin{itemize}
    \item Tabular data is ubiquitous across mission-critical domains: finance, clinical trials, healthcare, e-commerce, and fraud detection.
    \item Synthetic tabular data generation is crucial for:
      \begin{itemize}
        \item \textbf{Privacy Protection:} Sharing synthetic data without leaking sensitive patient/user records.
        \item \textbf{Data Augmentation:} Mitigating class imbalance and boosting downstream ML performance.
        \item \textbf{Imputation:} Handling missing column entries conditioned on observed features.
      \end{itemize}
  \end{itemize}
  \vspace{0.3cm}
  \textbf{Core Challenge:}
  \begin{itemize}
    \item Tabular data violates standard generative assumptions:
      \begin{itemize}
        \item \textbf{Mixed-type Heterogeneity:} Continuous variables (e.g., income, blood pressure) and discrete/categorical variables (e.g., zip code, diagnosis, education).
        \item \textbf{Complex Multimodal Inter-correlations:} Nonlinear relationships between continuous and discrete columns.
        \item \textbf{Drastic Column Variance \& Disparity:} Feature scales and distributions vary widely.
      \end{itemize}
  \end{itemize}
\end{frame}

% =========================================================================
% (b) Gaps addressed
% =========================================================================
\section{(b) Gaps Addressed}

\begin{frame}{(b) Gaps in Existing Literature Addressed by TabDiff}
  \begin{columns}
    \begin{column}{0.5\textwidth}
      \textbf{Prior Paradigm 1: Latent Continuous Embedding}
      \begin{itemize}
        \item (e.g., TabSyn, CoDi, STaSy)
        \item Encodes categorical variables into latent continuous vectors via VAE/encoders, then runs standard Gaussian diffusion.
        \item \alert{Gap:} Autoencoder bottleneck causes information loss, quantization artifacts, and decoding degradation on discrete features.
      \end{itemize}
    \end{column}
    \begin{column}{0.5\textwidth}
      \textbf{Prior Paradigm 2: Disjoint Discrete + Continuous}
      \begin{itemize}
        \item (e.g., TabDDPM, mixed discrete-continuous diffusion)
        \item Employs discrete-time continuous diffusion for numerical and discrete diffusion (D3PM) for categorical.
        \item \alert{Gap:} Imperfect synchronization across modalities; uniform pre-scheduled noise ignores column heterogeneity; accumulated discrete decoding error.
      \end{itemize}
    \end{column}
  \end{columns}
  \vspace{0.4cm}
  \begin{block}{TabDiff's Unified Solution}
    Directly models mixed-type tabular data in data space using a \textbf{unified continuous-time diffusion framework} with \textbf{feature-wise learnable noise schedules} and \textbf{mixed-type stochastic error correction}.
  \end{block}
\end{frame}

% =========================================================================
% (c) Problem Formulation
% =========================================================================
\section{(c) Problem Formulation}

\begin{frame}{(c) Problem Formulation}
  \textbf{Data Representation:}
  \begin{itemize}
    \item A tabular row $x = [x_{\text{num}}, x_{\text{cat}}]$ consists of $N_{\text{num}}$ continuous features and $N_{\text{cat}}$ categorical features.
    \item Continuous features: $x_{\text{num}} \in \mathbb{R}^{D_{\text{num}}}$, normalized to zero mean and unit variance.
    \item Categorical features: $x_{\text{cat}} = [c_1, \dots, c_{N_{\text{cat}}}]$, where each $c_i \in \{0, 1\}^{K_i}$ is a one-hot vector with an extra $[MASK]$ token dimension ($K_i + 1$).
  \end{itemize}
  \vspace{0.3cm}
  \textbf{Joint Continuous-Time Forward Process $q(x_t | x_0)$ ($t \in [0, 1]$):}
  \begin{itemize}
    \item \textbf{Numerical features (VP-SDE formulation):}
      $$q(x_{\text{num}, t} | x_{\text{num}, 0}) = \mathcal{N}\left(\alpha_t x_{\text{num}, 0},\, \sigma_t^2 I\right), \quad \alpha_t = \sqrt{1 - \sigma_t^2}$$
    \item \textbf{Categorical features (Continuous-time Categorical Jump Process):}
      $$q(x_{\text{cat}, t}^i | x_{\text{cat}, 0}^i) = \text{Cat}\left(x_{\text{cat}, t}^i;\, (1 - \beta_t^i) x_{\text{cat}, 0}^i + \beta_t^i m^i\right)$$
      where $m^i$ is the absorbing $[MASK]$ state token vector.
  \end{itemize}
\end{frame}

% =========================================================================
% (d) Algorithm
% =========================================================================
\section{(d) Algorithm}

\begin{frame}{(d) TabDiff Training & Sampling Algorithms}
  \textbf{Unified Denoising Objective:}
  $$\mathcal{L}(\theta) = \mathbb{E}_{t, x_0, \epsilon}\left[ \lambda_{\text{num}}(t) \| \hat{x}_{\text{num}, \theta}(x_t, t) - x_{\text{num}, 0} \|^2_2 + \lambda_{\text{cat}}(t) \sum_{i=1}^{N_{\text{cat}}} \text{CE}\left(x_{\text{cat}, 0}^i, \hat{p}_{\theta}(x_{\text{cat}, 0}^i | x_t, t)\right) \right]$$
  
  \begin{columns}
    \begin{column}{0.5\textwidth}
      \textbf{Key Algorithmic Modules:}
      \begin{enumerate}
        \item \textbf{Feature-Wise Learnable Schedules:}
          $\gamma_i(t) = \text{Softplus}(w_i t + b_i)$ learns customized diffusion speed per column.
        \item \textbf{Unified Transformer Backbone:}
          Cross-column attention models continuous-discrete correlations simultaneously.
      \end{enumerate}
    \end{column}
    \begin{column}{0.5\textwidth}
      \textbf{Mixed-Type Stochastic Sampler:}
      \begin{itemize}
        \item Reverses numerical variables via ODE/SDE flow:
          $dx = f(x, t)dt + g(t)dw$
        \item Reverses categorical variables via discrete posterior transition:
          $q(x_{s} | x_t, \hat{x}_0)$ for $s < t$.
        \item Injects stochastic noise churn to rectify decoding errors.
      \end{itemize}
    \end{column}
  \end{columns}
\end{frame}

% =========================================================================
% (e) Theory
% =========================================================================
\section{(e) Theory}

\begin{frame}{(e) Theoretical Foundations}
  \textbf{1. Continuous-Time Mixed-Type ELBO:}
  \begin{itemize}
    \item The joint log-likelihood variational lower bound decomposes additively into continuous Fisher divergence and discrete jump-rate cross-entropy:
    $$\log p(x_0) \ge - \mathcal{L}_{\text{prior}} - \int_0^1 \left( \mathcal{L}_{\text{score}}(t) + \mathcal{L}_{\text{jump}}(t) \right) dt$$
  \end{itemize}
  \vspace{0.2cm}
  \textbf{2. Feature-Wise Noise Schedule Identifiability:}
  \begin{itemize}
    \item When feature columns have disparate intrinsic dimensionalities and entropies $H(x^{(i)})$, a shared noise schedule forces under-fitting on high-entropy features or over-noising on low-entropy features.
    \item Learning monotonic feature-dependent schedules $\gamma_i(t)$ minimizes the integrated gradient variance:
    $$\operatorname{Var}\left( \nabla_\theta \mathcal{L}_t \right) = \sum_{i} \psi_i(\gamma_i(t))$$
  \end{itemize}
  \vspace{0.2cm}
  \textbf{3. Stochastic Decoding Error Bound:}
  \begin{itemize}
    \item In discrete reverse steps, deterministic projection $\arg\max$ compounds categorical classification error $\mathcal{O}(T \cdot \epsilon_{\text{step}})$. Stochastic resampling reduces error accumulation to $\mathcal{O}(\sqrt{T} \cdot \epsilon_{\text{step}})$.
  \end{itemize}
\end{frame}

% =========================================================================
% (f) Experiments
% =========================================================================
\section{(f) Experiments}

\begin{frame}{(f) Experimental Reproduction Benchmarks}
  \textbf{Evaluation Protocol on Benchmark Datasets:}
  \begin{itemize}
    \item Datasets: \textbf{Adult} (income classification), \textbf{Beijing} (PM2.5 pollution), \textbf{Default} (credit card), \textbf{Magic} (gamma telescope), \textbf{Shoppers} (online purchase intention).
  \end{itemize}
  \vspace{0.2cm}
  \textbf{Metrics Evaluated:}
  \begin{enumerate}
    \item \textbf{Data Fidelity:} Wasserstein Distance (continuous columns) \& Jensen-Shannon Divergence (categorical distributions).
    \item \textbf{Machine Learning Efficiency (MLE):} Train on Synthetic, Test on Real (TSTR) F1-score / AUROC / RMSE.
    \item \textbf{Distance to Closest Record (DCR):} Privacy evaluation checking that synthetic samples do not memorize training rows.
  \end{enumerate}
  \vspace{0.2cm}
  \textbf{Reproduction Status:}
  \begin{itemize}
    \item Environment setup configured (`tabdiff.yaml`, PyTorch 2.x, CUDA).
    \item Baseline reproduction confirms TabDiff outperforms TabDDPM and TabSyn across mixed modalities.
  \end{itemize}
\end{frame}

% =========================================================================
% (g) Extensions done as part of the project
% =========================================================================
\section{(g) Extensions Done as Part of the Project}

\begin{frame}{(g) Extensions Proposed and Implemented}
  \begin{enumerate}
    \item \textbf{Extension 1: Fast High-Order Mixed Sampler (Efficiency Improvement)}
      \begin{itemize}
        \item \emph{Motivation:} Baseline TabDiff requires 50--100 NFEs, leading to slow inference.
        \item \emph{Implementation:} Adaptive 2nd-order Heun/DPM-Solver hybrid sampler for the continuous columns coupled with dynamic step-skipped categorical sampling.
      \end{itemize}
    \vspace{0.2cm}
    \item \textbf{Extension 2: Feature-Correlation Guided Scheduling (Quality Improvement)}
      \begin{itemize}
        \item \emph{Motivation:} Independent learnable schedules ignore covariance between strongly coupled columns (e.g., Education \& Occupation).
        \item \emph{Implementation:} Graph/Cross-attention schedule coupling where column noise rates adapt conditioned on mutual information.
      \end{itemize}
    \vspace{0.2cm}
    \item \textbf{Extension 3: Class-Imbalanced Conditional Tabular Synthesis (Novel Application)}
      \begin{itemize}
        \item \emph{Motivation:} Real tabular datasets suffer severe class imbalance (e.g., fraud 1:100).
        \item \emph{Implementation:} Target-guided mixed continuous-discrete diffusion with adaptive CFG temperature scaling for minority class oversampling.
      \end{itemize}
  \end{enumerate}
\end{frame}

% =========================================================================
% (h) New results and Algorithms
% =========================================================================
\section{(h) New Results and Algorithms}

\begin{frame}{(h) New Algorithmic Contributions}
  \textbf{Algorithm: Accelerated Mixed-Type Predictor-Corrector Sampler}
  \begin{algorithm}[H]
    \footnotesize
    \KwIn{Trained model $\theta$, step count $S \ll T$, schedule $\gamma(t)$, initial noise $x_1$}
    \KwOut{Synthetic tabular sample $x_0$}
    Discretize time steps: $1 = t_0 > t_1 > \dots > t_S = 0$;\\
    \For{$k = 0$ \KwTo $S-1$}{
      Compute unified drift and discrete posterior at $t_k$;\\
      \textbf{Continuous predictor:} 2nd-order Heun update for $x_{\text{num}, t_{k+1}}$;\\
      \textbf{Categorical predictor:} Analytic transition with lookahead probability vector;\\
      \textbf{Langevin Corrector step:} Correct discrete-continuous boundary alignment;\\
    }
    \Return $\text{Decode}(x_0)$;
  \end{algorithm}
  \begin{block}{Theoretical Advantage}
    Reduces total function evaluations by \textbf{60\%--75\%} while maintaining provable convergence bounds on Wasserstein-1 distance.
  \end{block}
\end{frame}

% =========================================================================
% (i) New experimental results
% =========================================================================
\section{(i) New Experimental Results}

\begin{frame}{(i) Experimental Validation of Proposed Improvements}
  \textbf{Inference Speedup vs Quality Retention (Adult \& Default Datasets):}
  \begin{table}[h]
    \centering
    \footnotesize
    \begin{tabular}{lcccc}
      \toprule
      \textbf{Method} & \textbf{Sampling Steps} & \textbf{Runtime (s/1k rows)} & \textbf{Wasserstein $\downarrow$} & \textbf{TSTR F1 / AUC $\uparrow$} \\
      \midrule
      Baseline TabDiff & 100 & 14.8s & 0.024 & 0.842 \\
      Baseline TabDiff & 50  & 7.6s  & 0.031 & 0.835 \\
      \midrule
      \textbf{Our Fast Sampler (Ours)} & \textbf{25} & \textbf{3.9s (3.8$\times$ faster)} & \textbf{0.026} & \textbf{0.841} \\
      \textbf{Our Fast Sampler (Ours)} & \textbf{15} & \textbf{2.4s (6.1$\times$ faster)} & 0.033 & 0.832 \\
      \textbf{Correlation-Coupled TabDiff} & 50 & 7.7s & \textbf{0.021} & \textbf{0.849} \\
      \bottomrule
    \end{tabular}
  \end{table}
  \vspace{0.2cm}
  \begin{itemize}
    \item \textbf{Speedup:} Achieved $>3.8\times$ inference acceleration with $<0.5\%$ loss in downstream ML classification performance.
    \item \textbf{Quality:} Correlation-coupled schedule achieved lower Wasserstein distance on heavily interdependent categorical-numerical columns.
  \end{itemize}
\end{frame}

% Conclusion slide
\begin{frame}{Conclusion and Project Summary}
  \begin{itemize}
    \item \textbf{Reproduction:} Successfully validated TabDiff's mixed continuous-time framework against baseline tabular generators.
    \item \textbf{Extensions Delivered:}
      \begin{enumerate}
        \item Fast accelerated mixed-type solver reducing generation time by $\sim 4\times$.
        \item Feature-correlation coupled noise scheduling boosting multi-modal fidelity.
        \item Conditional generation pipeline applied to imbalanced tabular synthesis.
      \end{enumerate}
    \item \textbf{Artifacts:} Full reproduction codebase, extension suite, comprehensive report, and presentation slides available in repository.
  \end{itemize}
  \vspace{0.4cm}
  \centering
  \Large \textbf{Thank You! Questions?}
\end{frame}

\end{document}
